\section{Data and Methodology}
\label{sec:method}
The design has three stages. We reproduce the industry benchmark, measure scalar restriction and uncertainty-estimation losses in a known model, and compare feasible allocations on historical returns. In the historical comparison, scalar and heterogeneous adjustments use the same return forecasts, covariance estimates, scalar-selected penalty parameter and trading assumptions. The allocation input that changes is relative uncertainty: it is one for every asset in the scalar rule and estimated separately for each asset in the heterogeneous rule. This change can affect composition, exposure and turnover. Table~\ref{tab:design} summarizes the settings.

\subsection{Data, timing and common risk inputs}
We use monthly value-weighted returns on the 49 industry portfolios \citep{frenchdata} and the monthly risk-free series from the Fama/French Research Data Factors file \citep{frenchfactors}. Both datasets are from the Kenneth R.\ French Data Library. The downloaded files extend from July 1926 to August 2026; evaluation ends in December 2025. The benchmark covers 1942--2018. The extension reports 1970--2018 and 2019--2025 separately. All dates in performance tables denote return realization months. Information available at the end of month $t$ determines positions held during month $t+1$.

Let $r_{i,t}$ be industry $i$'s total return, $r_{f,t}$ the risk-free return and $r^e_{i,t}=r_{i,t}-r_{f,t}$ its excess return. Returns are converted from percentages to decimals; missing-value codes at or below $-99.99$ in the source percentage units are treated as unavailable. At formation date $t$, the benchmark universe contains industries with 60 consecutive observed monthly returns through $t$, positive sample variance and a complete 12-month momentum history. At least five industries are required. Let $\widehat\Sigma_t$ denote the sample covariance of excess returns over these 60 months, computed with denominator 59, and let $V_t=\diag(\widehat\Sigma_t)$. The preliminary risk estimate is
\begin{equation}
Q_t=0.95\widehat\Sigma_t+0.05V_t.
\label{eq:risk}
\end{equation}
The diagonal is unchanged. We apply this 5\% correlation shrinkage before the selected EPO adjustment. The matrix $Q_t$ measures return covariance risk, while the subsequent uncertainty adjustment modifies the allocation to account for signal uncertainty.

We first compute each industry's compounded total return over months $t-11$ through $t$. We subtract its cross-sectional mean and scale the positive and negative scores so that they sum to $+1$ and $-1$, respectively. Denote the resulting score by $z_{i,t}$. The benchmark uses this normalized momentum score directly as its signal. The extension uses the score as a predictor in a rolling regression to estimate next-month expected excess returns. The score is formed on the benchmark universe before applying the extension's stricter forecast eligibility rule. Appendix~\ref{app:data} specifies the data vintage and benchmark sample accounting.
\input{generated/table_design}

\subsection{Replication design}
The benchmark compares raw mean--variance optimization, its covariance-corrected baseline and EPO. Their unscaled portfolio directions are
\begin{equation}
\begin{aligned}
d_t^{\mathrm{raw}}&=\widehat\Sigma_t^{-1}z_t,\\
d_t^{Q}&=Q_t^{-1}z_t,\\
d_t^{\mathrm{EPO}}(w)&=[(1-w)Q_t+wV_t]^{-1}z_t,
\end{aligned}
\label{eq:benchmark-directions}
\end{equation}
where $w\in\{0,0.25,0.50,0.75,0.90,0.95,0.99\}$. Let $j$ index these benchmark rules, including each fixed-$w$ EPO candidate. For a nonzero direction $d_t^{(j)}$, the actual position vector is
\begin{equation}
x_t^{(j)}=
\frac{v_{\mathrm{target}}}
{\sqrt{12\,d_t^{(j)\prime}Q_t d_t^{(j)}}}\,d_t^{(j)},
\qquad v_{\mathrm{target}}=0.10.
\label{eq:benchmark-scaling}
\end{equation}
A zero direction remains a zero position. Because $Q_t$ is a monthly covariance matrix, the factor 12 converts predicted variance to annual units. All benchmark rules therefore use the same risk matrix for their 10\% predicted-volatility target, including the raw optimizer. Because positions are scaled to a common ex ante volatility, only the direction $d_t^{\mathrm{EPO}}(w)$ matters, and the zero-anchor allocation of Section~\ref{sec:framework} coincides with the simple EPO of \citet{pedersen2021}. The empirical extension compares direct fixed-risk-aversion allocations with a common volatility target and gross cap, as specified below. The reported benchmark returns are gross excess returns. In January of each year we choose the candidate with the largest annualized gross Sharpe ratio over the preceding 60 complete calendar months. Ties within $10^{-12}$ select the larger $w$. The selected value remains fixed for that year, while signals and positions update monthly. A common valid-month intersection is used for the reported comparison.

We also reproduce the principal-component mechanism. Each month we diagonalize the correlation matrix of each portfolio's risk matrix ($\widehat\Sigma_t$, $Q_t$ or the selected EPO matrix), order its eigenvalues from smallest to largest and place components into five groups; the three correlation matrices share eigenvectors, so the groups coincide. For each group we compare predicted and realized risk, signal-implied and realized performance, and the share of predicted portfolio variance. These diagnostics distinguish the contribution of small-eigenvalue directions from the aggregate performance comparison. Appendix~\ref{app:math} defines the component statistics exactly.

\subsection{Simulation design}
\label{sec:simulation-method}
The first experiment implements equations~\eqref{eq:model}--\eqref{eq:loss} with $a=0$, $n=49$, $\gamma=3$, $\tau=0.04$ and $\lambda=0.10$. Every asset has monthly volatility 0.05 and every pair has correlation 0.30. For each target cross-sectional coefficient of variation $c\in\{0,0.5,1\}$, we construct a positive uncertainty profile with mean one. For $c>0$, draw independent $v_i\sim N(0,1)$ and set
\begin{equation}
\ell_c=\sqrt{\log(1+c^2)},
\qquad
h_i=\frac{\exp(\ell_c v_i-\ell_c^2/2)}
{n^{-1}\sum_{j=1}^{n}\exp(\ell_c v_j-\ell_c^2/2)}.
\label{eq:simulation-profile}
\end{equation}
For $c=0$, set every $h_i=1$. The lognormal population coefficient of variation is $c$; the realized cross-sectional coefficient in a finite profile can differ, so both actual and target dispersions are reported.

We draw one profile for each value of $c$ before generating expected returns or signal errors, and hold it fixed across $K=1{,}000$ draws. Each draw generates $\mu$, $\varepsilon$ and $s=\mu+\varepsilon$ from equation~\eqref{eq:model} with $a=0$ and the profile determining $\Lambda=\lambda\diag(\sigma_i^2h_i)$. The same underlying standard-normal innovations are reused across profiles. For every signal draw we compute the conditional mean, the heterogeneous allocation, the best scalar allocation and the mean-scalar allocation. The scalar search evaluates 2,001 equally spaced values of $w$ on $[0,1]$, refines the best grid neighborhood by golden-section search, and compares the endpoints. The value $w=1$ is the zero-position limit.

For draw $k$, write $s_k$ for the signal, $x_k^H$ for its heterogeneous allocation, $L_k=L(s_k)$ for compression loss and $U_k^H=\bar U(x_k^H\mid s_k)$ for heterogeneous utility. The reported mean loss and proportional loss are
\begin{equation}
\overline L=\frac{1}{K}\sum_{k=1}^{K}L_k,
\qquad
\overline U^H=\frac{1}{K}\sum_{k=1}^{K}U_k^H,
\qquad
\text{Loss share}=\frac{\overline L}{\overline U^H}.
\label{eq:simulation-summary}
\end{equation}
We express the loss share as a percentage by multiplying by 100. It is a ratio of sample means; individual-draw utility ratios are not averaged. We also report mean additional choice loss. The model and its utility are monthly. The main text and figures report utility shares. In Appendix~\ref{app:simulation}, we multiply utility levels, losses, Monte Carlo standard errors and interval endpoints by $12\times10^4$ to express them in annual basis points. This is arithmetic annualization of monthly quadratic utility; it does not simulate an additional twelve-month holding process or compound returns. Loss shares are unchanged. These conditional model losses measure the cost of a restriction with known inputs, whereas the historical certainty equivalents summarize realized net returns; the common time unit does not make them interchangeable economic gains. Forecast errors and bootstrap standard errors retain monthly units because they refer to next-month predictions. Within each fixed profile, draws are divided into thirds by the exposure-alignment statistic in equation~\eqref{eq:alignment}. Appendix~\ref{app:simulation} specifies the seed, conditional Monte Carlo intervals and calibration design.

The second experiment measures the cost of estimating relative uncertainty. We retain the true $h$, signals and conditional means from the first experiment, and fix the common scalar at $\bar\kappa=\lambda/\tau=2.5$. For draw $k$, generate independent $\zeta_{i,k}\sim N(0,1)$, independent of the true signal and its error. At estimation-error levels $\delta\in\{0,\delta_{\mathrm{cal}},0.5,1\}$, with $\delta_{\mathrm{cal}}$ defined in Section~\ref{sec:forecast-method}, construct
\begin{equation}
\widetilde h_{i,k}(\delta)=h_i\exp(\delta\zeta_{i,k}-\delta^2/2),\qquad
\widehat h_{i,k}(\delta)=\frac{\widetilde h_{i,k}(\delta)}{n^{-1}\sum_j\widetilde h_{j,k}(\delta)}.
\label{eq:controlled-h-error}
\end{equation}
The same standard-normal draws are used across error levels and true-uncertainty settings. The estimated matrix is $\widehat\Lambda_k(\delta)=\lambda\diag(\sigma_i^2\widehat h_{i,k}(\delta))$. We form $\widehat x^H$ from this matrix and the original signal, then compute $L_{\bar\kappa}$, $P$ and their difference using the true conditional objective. The experiment isolates errors in relative uncertainty while keeping its common scale known. The parameter $\delta$ specifies controlled measurement error. The additional level $\delta_{\mathrm{cal}}$ matches the estimator's log-error dispersion in synthetic calibration, as defined in Section~\ref{sec:forecast-method}. At $\delta=0$, $P=0$; with uniform true uncertainty, $L_{\bar\kappa}=0$. We verify equation~\eqref{eq:estimation-tradeoff} for every draw and report mean losses, their shares of mean oracle utility and the fraction with $P<L_{\bar\kappa}$. Appendix~\ref{app:simulation} gives numerical results and Monte Carlo intervals.

\subsection{Forecasts and relative uncertainty}
\label{sec:forecast-method}
For the historical extension we estimate an intercept and slope separately for each industry by ordinary least squares. At formation date $t$, the regression contains exactly 120 pairs $(z_{i,j},r^e_{i,j+1})$ for $j=t-120,\ldots,t-1$. Writing its coefficients as $\widehat\alpha_{i,t}$ and $\widehat\beta_{i,t}$, the next-month forecast is
\begin{equation}
\widehat s_{i,t}=\widehat\alpha_{i,t}+\widehat\beta_{i,t}z_{i,t}.
\label{eq:forecast}
\end{equation}
The latest response is observed in month $t$; the return in $t+1$ is never used to fit its forecast. An industry must have all 120 consecutive pairs, a current score and an admissible risk history. Entry occurs at its first forecast-ready month. We later require 60 consecutive forecast-ready months as a sensitivity exercise.

We use a moving-block pairs bootstrap, following the dependent-data resampling approach of \citet{kunsch1989}. A bootstrap sample concatenates uniformly sampled overlapping blocks of 12 consecutive training pairs and truncates the concatenation to 120 pairs. Blocks do not wrap around the sample. The same sampled time indices apply to all industries. Let $B=200$ be the number of bootstrap samples, and let $\widehat s_{i,t}^{*(k)}$ be the forecast from the regression refitted on sample $k$, evaluated at the fixed current score $z_{i,t}$. Define the bootstrap mean and forecast variance by
\begin{equation}
\overline s_{i,t}^{*}=\frac{1}{B}\sum_{k=1}^{B}\widehat s_{i,t}^{*(k)},
\qquad
u_{i,t}=\frac{1}{B-1}\sum_{k=1}^{B}
\left(\widehat s_{i,t}^{*(k)}-\overline s_{i,t}^{*}\right)^2.
\label{eq:bootstrap-variance}
\end{equation}
Thus $u_{i,t}$ measures variation in the estimated conditional mean across resampled training histories. Future return innovations are not added to these forecasts. Resampled efficiency averages portfolios optimized under simulated investment inputs \citep[Chapter~6]{michaud2008}. Here resampling is used only to estimate forecast variance: the original regression supplies the point forecast, and one portfolio is formed from that forecast and the estimated uncertainty profile at each date.

Let $\mathcal I_t$ be the eligible-industry set and $N_t=|\mathcal I_t|$ its size. For each industry, divide $u_{i,t}$ by its corresponding 60-month return variance $\widehat\sigma_{i,t}^2=(V_t)_{ii}$. Let $\operatorname{Winsor}_t$ clip these ratios to their cross-sectional 5th and 95th percentiles within $\mathcal I_t$, using linearly interpolated empirical quantiles. The relative uncertainty profile is
\begin{equation}
\begin{aligned}
\eta_{i,t}&=\frac{u_{i,t}}{\widehat\sigma_{i,t}^{2}},
&\widetilde\eta_{i,t}&=\max\left\{10^{-8},\operatorname{Winsor}_t(\eta_{i,t})\right\},\\
\widehat h_{i,t}&=\frac{\widetilde\eta_{i,t}}
{N_t^{-1}\sum_{j\in\mathcal I_t}\widetilde\eta_{j,t}},
&\frac{1}{N_t}\sum_{i\in\mathcal I_t}\widehat h_{i,t}&=1.
\end{aligned}
\label{eq:relative-uncertainty}
\end{equation}
Scaling by return variance expresses estimation uncertainty relative to asset risk. Winsorization limits extreme ratios, and normalization fixes the cross-sectional mean of $\widehat h$ at one. Portfolio exposure and risk remain outcomes of the allocation rule. If all forecast variances are zero, set all $\widehat h_{i,t}=1$. The covariance matrix is recomputed on the actual eligible universe and the uncertainty profile is normalized on that same universe. Bootstrap seeds depend on formation month, so recomputing a date produces the same forecasts.

The quantity $u_{i,t}$ estimates sampling variance of the fitted conditional mean. The normalized $\widehat h_{i,t}$ supplies relative uncertainty to the allocation, and the selected scalar controls the common penalty. Calibration evaluates this estimator against a known conditional mean, using the same regression window and resampling procedure.

Before portfolio evaluation, we run the identical estimator on 1,000 synthetic histories with a known conditional-mean regression. Persistent predictors have autoregressive coefficient 0.92; residual volatilities span 0.04 to 0.08 and residual cross-industry correlation is 0.30. We study both independent predictor innovations and a feedback specification in which a return innovation enters the next predictor. Each history supplies 120 training pairs and a current forecast. We compare estimated variance with squared conditional-mean estimation error, evaluate interval coverage, and assess relative-uncertainty rankings. The full process, including the feedback coefficient and the appropriate interpretation of its analytical comparison, is specified in Appendix~\ref{app:simulation}.

The zero-feedback design also provides the log-error dispersion used in the controlled experiment of Section~\ref{sec:simulation-method}. For history $k$, let $h^\star_{i,k}$ be the conditional OLS forecast-error variance divided by population return variance, normalized to cross-sectional mean one, and set $e_{i,k}=\log(\widehat h_{i,k}/h^\star_{i,k})$. With $K=1{,}000$ histories, $n=49$ industries and $\bar e_k=n^{-1}\sum_i e_{i,k}$, define
\begin{equation}
\delta_{\mathrm{cal}}=\left[\frac{1}{K}\sum_{k=1}^K\frac{1}{n-1}\sum_{i=1}^n(e_{i,k}-\bar e_k)^2\right]^{1/2}.
\label{eq:calibrated-error}
\end{equation}
Within-history centering removes common normalization factors. This matches the second moment of cross-sectional log error to the controlled Gaussian error scale. We evaluate the controlled experiment at this measured value, retaining its independent Gaussian errors and fixed true uncertainty settings.

\subsection{Allocations, selection and position sizing}
Let $\widehat H_t=\diag(\widehat\sigma_{i,t}^2\widehat h_{i,t})$ on the eligible universe. With fixed risk aversion $\gamma=3$ and $\kappa_t=w_t/(1-w_t)$, the direct-allocation portfolios are
\begin{equation}
x_t^S=\frac{1}{\gamma}(Q_t+\kappa_t V_t)^{-1}\widehat s_t,
\qquad
x_t^H=\frac{1}{\gamma}(Q_t+\kappa_t\widehat H_t)^{-1}\widehat s_t.
\label{eq:empirical}
\end{equation}
Equation~\eqref{eq:empirical} implements the EPO allocation using the same estimated mean, risk matrix and scalar level, with relative uncertainty estimated by bootstrap for the heterogeneous rule. Here $\kappa_t$ is an effective penalty parameter selected from past portfolio performance.

The extension selects one common penalty each January using the preceding 60 complete months of scalar candidate performance. Let $Y$ denote the holding year and $t_Y$ the preceding December. For a fixed grid value $w$, let $y_m^S(w)$ be the scalar candidate's net excess return realized in month $m$, including trading and short costs. Define $\mathcal W_Y$ as the subset of the seven-value grid whose candidates have not previously become insolvent and have all 60 required monthly returns available. With $\CE$ denoting the annualized sample certainty equivalent defined in equation~\eqref{eq:ce}, the selected value is
\begin{equation}
w_Y\in\operatorname*{arg\,max}_{w\in\mathcal W_Y}
\CE\!\left(\{y_m^S(w)\}_{m=t_Y-59}^{t_Y}\right).
\label{eq:common-selection}
\end{equation}
Numerical ties within $10^{-12}$ select the larger $w$. Both portfolios use this $w_Y$ for all return months in year $Y$; in equation~\eqref{eq:empirical}, $w_t=w_Y$ when the position formed at $t$ earns its return in that year. Candidate accounts begin from cash in January 1960 and maintain their own trading histories. A candidate whose total wealth factor becomes nonpositive is permanently excluded from selection. The chosen portfolios begin from cash in January 1970; their positions and wealth evolve continuously through December 2025, including the reporting boundary in 2019. Thus the later period does not receive a new initial account. The direct-allocation benchmark has fixed risk aversion and no gross-exposure cap.

This design measures the effect of asset-specific uncertainty at a common scalar-selected penalty. The main economic comparison applies a common position-sizing rule; direct fixed-$\gamma$ allocations provide the benchmark for assessing the effects of scale. For $j\in\{S,H\}$, set
\begin{equation}
\begin{aligned}
\widetilde x_t^j&=\frac{0.10}{\sqrt{12\,x_t^{j\prime}Q_tx_t^j}}\,x_t^j,\\
x_t^{j,\mathrm{target}}&=\widetilde x_t^j\min\left\{1,\frac{3}{\sum_i|\widetilde x_{i,t}^j|}\right\}.
\end{aligned}
\label{eq:deployment}
\end{equation}
A zero portfolio remains zero. This sets predicted annual volatility to 10\% before applying a gross-exposure cap of three. The cap can reduce predicted risk below the target. This comparison retains the fixed-$\gamma$ rule's selected $w_t$ sequence; each strategy carries its own holdings, drift and costs. Risk aversion remains three when evaluating utility under either position-sizing rule.

To describe portfolio composition, define gross exposure $g^j_t=\sum_i|x^j_{i,t}|$ for strategy $j\in\{S,H\}$. Gross-normalized composition distance is $\sum_i|x^H_{i,t}/g^H_t-x^S_{i,t}/g^S_t|$. Risk-direction correlation is
\begin{equation}
\rho_{Q,t}=\frac{x_t^{S\prime}Q_t x_t^H}
{\sqrt{x_t^{S\prime}Q_t x_t^S}\sqrt{x_t^{H\prime}Q_t x_t^H}}.
\label{eq:direction-correlation}
\end{equation}
A positive rescaling alone gives distance zero and correlation one. These measures distinguish changes in composition from changes in overall exposure.

\subsection{Trading and evaluation}
Let $x^-_{i,t}$ denote the position immediately before rebalancing at date $t$, expressed relative to current wealth. Full traded notional is $T_t=\sum_i|x_{i,t}-x^-_{i,t}|$, including liquidation of industries leaving the universe. Short notional is $B_t=\sum_i\max(-x_{i,t},0)$. With trading rate $c=0.001$ per dollar traded and annual short fee $b=0.005$, net excess return is
\begin{equation}
y_{t+1}=x_t'r^e_{t+1}-cT_t-\frac{b}{12}B_t.
\label{eq:net}
\end{equation}
The residual cash position $1-\sum_i x_{i,t}$ earns or pays the risk-free rate. Total wealth grows by $G_{t+1}=1+r_{f,t+1}+y_{t+1}$. With fees charged at month end, the next pretrade risky position is $x^-_{i,t+1}=x_{i,t}(1+r_{i,t+1})/G_{t+1}$. Turnover uses the full absolute sum, without dividing by two. Missing realized returns on held positions cause an explicit failure rather than an imputed portfolio return.

For $M$ monthly net excess returns, let $\bar y$ and $s_y^2$ be their sample mean and sample variance, the latter using denominator $M-1$. The annualized certainty equivalent is
\begin{equation}
\CE(y)=12\bar y-\frac{\gamma}{2}\,12s_y^2,
\qquad \Delta\CE=\CE(y^H)-\CE(y^S).
\label{eq:ce}
\end{equation}
Annualized Sharpe ratio is $\sqrt{12}\bar y/s_y$ and volatility is $\sqrt{12}s_y$. We also report monthly expected shortfall at 95\%, maximum drawdown from total net wealth, full turnover and gross exposure. Expected shortfall averages the worst 5\% of net excess returns, with fractional weight at the empirical boundary; it is reported as a positive loss. Drawdown includes initial wealth one within each reporting period.

For each position-sizing rule and period, uncertainty intervals for $\Delta\CE$ use 1,000 paired moving-block bootstrap samples of realized monthly net returns. We sample overlapping blocks of 12 consecutive months uniformly, concatenate them and truncate to the original period length. All strategies use the same sampled dates. Each bootstrap sample recomputes means, sample variances and the CE difference; the 2.5th and 97.5th percentiles form the interval. The procedure resamples the returns generated by the fitted forecasts and selection rules. Blocks of 6 and 24 months assess sensitivity under both position-sizing rules.

For forecast evaluation, let $\mathcal O$ contain the eligible industry--formation-month pairs $(i,t)$ whose next-month returns fall within the reported evaluation period. The historical-mean benchmark uses the same 120 completed responses as the regression. Define
\begin{equation}
\begin{aligned}
\overline r^{\,\mathrm{hist}}_{i,t}
&=\frac{1}{120}\sum_{j=t-119}^{t}r^e_{i,j},\\
R_{\mathrm{OOS}}^2
&=1-
\frac{\sum_{(i,t)\in\mathcal O}
\left(r^e_{i,t+1}-\widehat s_{i,t}\right)^2}
{\sum_{(i,t)\in\mathcal O}
\left(r^e_{i,t+1}-\overline r^{\,\mathrm{hist}}_{i,t}\right)^2}.
\end{aligned}
\label{eq:forecast-r2}
\end{equation}
Here OOS denotes out-of-sample prediction. The ratio pools squared errors over all eligible industry-month observations. A negative value means the fitted forecasts have higher squared error than the historical-mean benchmark. We also calculate the monthly information coefficient as the cross-sectional Spearman correlation between forecasts and next-month excess returns. Intervals for these diagnostics resample whole months in blocks, preserving contemporaneous industry dependence.

To evaluate uncertainty-to-asset assignment, we compare the original $\widehat h_t$ with 100 random assignments to industries. A fixed random ranking determines each assignment, restricted to the eligible industries at each date. This preserves the monthly distribution of uncertainty estimates and keeps the mapping fixed while eligibility is unchanged. Forecasts, risk estimates and selected penalties are common; each portfolio generates its own holdings and costs under both position-sizing rules. Let $y^O$ denote the original strategy's net-return series and $y^{(b)}$ that of random assignment $b$. Define
\begin{equation}
D=\CE(y^O)-\frac{1}{B_{\mathrm{A}}}\sum_{b=1}^{B_{\mathrm{A}}}\CE(y^{(b)}),
\qquad B_{\mathrm{A}}=100.
\label{eq:assignment-advantage}
\end{equation}
The benchmark averages the utilities of the separately implemented strategies. We calculate an interval for $D$ by applying the same paired moving-block bootstrap to all 101 return series together, keeping the assignments fixed. This is a time-series comparison with a specified random-assignment benchmark. The percentage of assignments with lower utility than the original is also reported as a descriptive rank. Appendix~\ref{app:label-placebo} specifies the mapping and resampling details.

Two trading-cost scenarios set $c$ to zero or 0.0025 under both position-sizing rules, retaining the selected $w$ sequence and recomputing each strategy's position drift, turnover and costs. The short fee remains unchanged. A supporting 60-ready-month admission exercise uses fixed risk aversion and reruns candidate histories and selection. These comparisons use the same excess-return and utility definitions as the baseline.
